% ch3-part1.tex — Chapter 3 intro + Sprints 1–3
% Compile with XeLaTeX / LuaLaTeX + fontspec (unicode arrows, checkmarks, accents kept as-is).

\chapter{Development Sprints}

This chapter narrates the five development sprints that built the platform, in the dependency
order fixed by the sprint plan in Section~2.3. Each sprint opens with its goal and backlog,
analyses the requirements per module, presents the UML design alongside the features it
describes, reports the implementation with screenshots, validates the deliverables through
testing, and closes with a review, retrospective, and cumulative class diagram showing how the
domain model grew. The AI sprint, being the technical climax, is told separately in Chapter~4.

%% ====================================================================
%% SPRINT 1
%% ====================================================================

\clearpage
\section{Sprint 1: Foundations \& Authentication}

\subsection{Sprint goal}

Stand up everything the rest of the platform stands on: a layered NestJS backend with a Prisma
multi-file schema, and the identity core --- authentication (login, token refresh, password
reset), role-based access control, and administration of users and teams. Nothing in later
sprints can be built or even reached until a caller can be identified and authorized, which is
why this work comes first.

\subsection{Sprint backlog}

The user stories for this sprint, repeated from the product backlog in Section~2.2 so the
sprint reads self-contained:

\begin{longtable}{@{}l>{\raggedright\arraybackslash}p{0.50\textwidth}lll@{}}
\caption{Sprint 1 backlog.}\label{tab:s1-backlog-detail} \\
\tableheaderrow
\thcell{US-ID} & \thcell{Story} & \thcell{SP} & \thcell{Est.} & \thcell{Priority} \\
\midrule
\endfirsthead
\caption[]{Sprint 1 backlog (continued).} \\
\tableheaderrow
\thcell{US-ID} & \thcell{Story} & \thcell{SP} & \thcell{Est.} & \thcell{Priority} \\
\midrule
\endhead
\bottomrule
\endfoot
\rowcolor{rowalt}
US-S1-01 & As a developer, I want a layered NestJS backend with a Prisma multi-file schema and migrations, so that every module shares one structure. & 8 & 4d & Must \\
US-S1-02 & As a user, I want to log in with email + password and receive a JWT, so that I can access the platform. & 3 & 1.5d & Must \\
\rowcolor{rowalt}
US-S1-03 & As a user, I want my session refreshed automatically on expiry, so that I stay logged in without re-authenticating. & 3 & 1d & Must \\
US-S1-04 & As a user, I want to reset a forgotten password via an emailed code, so that I can regain access. & 3 & 2d & Should \\
\rowcolor{rowalt}
US-S1-05 & As the platform, I want every protected endpoint gated by a role-permission guard, so that a user only reaches what their role allows. & 8 & 3.5d & Must \\
US-S1-06 & As an admin/HR, I want to provision user accounts directly with one or more roles, so that staff can be onboarded immediately without going through self-registration's pending-approval step. & 5 & 2.5d & Must \\
\rowcolor{rowalt}
US-S1-07 & As an admin, I want to list, search (accent-insensitive), and filter users, so that I can manage the directory. & 3 & 1.5d & Should \\
US-S1-08 & As an admin, I want to edit and soft-delete (deactivate) a user, so that I manage the account lifecycle. & 3 & 1.5d & Must \\
\rowcolor{rowalt}
US-S1-09 & As a user, I want to update my own profile and change my password, so that I manage my account. & 2 & 1d & Should \\
US-S1-10 & As an admin, I want to create teams and assign members and a manager, so that staff are grouped. & 3 & 1.5d & Should \\
\midrule
& \textbf{Sprint 1 subtotal} & \textbf{41} & \textbf{3w} & \\
\end{longtable}

\subsection{Requirements analysis}

The sprint delivers two modules. The features, actors, and use cases are collected in
Figure~\ref{fig:3-1}.

{\renewcommand{\figmaxheight}{0.75\textheight}%
\reportfigure[\textwidth]{figures/fig_3_1}{Sprint 1 use-case diagram.}{3-1}}

Three actor archetypes appear. A \textbf{Visitor} (unauthenticated) can self-register (landing in a
\texttt{PendingApproval} state until promoted by an admin) or reset a password. A \textbf{User} (any
authenticated employee) manages their own profile and credentials. An \textbf{Admin / HR} provisions
and administers accounts, roles, and teams.

\subsubsection{Module A --- Authentication \& RBAC}

\textbf{Functional description.} This module is the cross-cutting security layer every other
module depends on. It proves \emph{who} is calling (identity) and decides \emph{what} they may
do (authorization). Users authenticate with email and password; the API also accepts phone-based
login, though the current client uses email only. On success the server returns a JWT
access/refresh pair. A three-step password-reset flow (request a 5-digit
code, verify it, set a new password) uses a single live code per user. Refresh tokens are the
only server-side authentication state --- the token string itself is the primary key of the
\texttt{RefreshToken} table --- which lets logout revoke a session and lets the refresh endpoint
reject a token that was never issued.

\textbf{Actors.} A Visitor initiates login or password reset. A User triggers token refresh and
logout. The platform itself (via the guard) enforces permissions on every protected route.

\textbf{Business rules.}
\begin{itemize}
  \item A user may hold \textbf{multiple roles simultaneously}; roles are modeled as rows in a
    \texttt{Role} table rather than a single enum column, and login flattens them into the JWT
    payload.
  \item Authorization is a \textbf{static, compile-time role-to-permission map}: approximately
    120 fine-grained permissions mapped to 31 role types in a shared constant
    (\texttt{PERMISSIONS\_FOR\_ROLE}).
  \item A single \texttt{HasPermissionGuard}, applied on 139 routes across 18 controllers, checks
    that at least one of the caller's roles carries at least one of the route's required
    permissions (\textbf{OR} across permissions, \textbf{OR} across roles).
  \item Finer ownership distinctions (\texttt{*.own} vs \texttt{*.any}) are resolved inside each
    service, where the data needed to make that call lives.
  \item Only one password-reset code may be active per user at a time.
  \item Self-registration lands in a \texttt{PendingApproval} state with minimal permissions (a
    least-privilege default); only a privileged role can promote the account.
\end{itemize}

\textbf{Login sequence.} Figure~\ref{fig:3-2} shows the login path: the service fetches the user
by email or phone (matching only active accounts), compares the password with bcrypt, and on
success mints the token pair and persists the refresh row.

\reportfigure{figures/fig_3_2}{Login sequence diagram.}{3-2}

\textbf{Protected-request sequence.} Figure~\ref{fig:3-3} shows how every subsequent request is
authorized: the guard reads the route's \texttt{@Permissions} metadata, verifies the JWT
signature and expiry, and attaches the decoded payload to the request. If no required permission
intersects the caller's effective set, the guard returns 403.

\reportsideways{figures/fig_3_3}{Protected-request RBAC sequence.}{3-3}

\textbf{Class diagram.} Figure~\ref{fig:3-4} shows the identity core: \texttt{User} holds
credentials and the active flag; \texttt{Role} attaches one or more \texttt{UserType} values to a
user; \texttt{RefreshToken} records live sessions; and \texttt{ResetPasswordCode} holds the single
outstanding reset code.

\reportfigure[0.90\textwidth]{figures/fig_3_4}{Auth \& RBAC class diagram.}{3-4}

\subsubsection{Module B --- Users \& Teams}

\textbf{Functional description.} On top of the identity core, this module runs the
administrative user lifecycle --- create, read, update, soft-delete, profile self-service,
password change, and bulk email --- and groups users into teams. Account provisioning is a
privileged action: an admin or HR user creates accounts that are immediately active with the
requested roles. This is the counterpart to self-registration in Module~A, which produces an
inert, pending account.

\textbf{Actors.} An \textbf{Admin / HR} provisions and manages user accounts. A \textbf{User}
manages their own profile and password.

\textbf{Business rules.}
\begin{itemize}
  \item Authorization is \textbf{two-tiered}: the coarse \texttt{HasPermissionGuard} gates the
    route, then a fine-grained, data-scoped check runs inside the service. This pattern is
    reused across the entire platform.
  \item An admin can only grant roles they are themselves allowed to manage
    (\texttt{canRolesManageRoles}), preventing privilege escalation (e.g., an HR manager cannot
    create a CEO account).
  \item User names are stored twice: the original \texttt{name} and a denormalized,
    accent-stripped \texttt{unaccentedName} for accent-insensitive search.
  \item User deletion is a \textbf{soft delete}: the \texttt{isActive} flag is set to
    \texttt{false} rather than the row removed, preserving history and foreign-key integrity.
  \item Teams group users through a \texttt{UserTeam} join table; the \texttt{isManager} flag on
    that join encodes who manages whom for later authorization decisions.
\end{itemize}

\textbf{Create-user sequence.} Figure~\ref{fig:3-5} shows the admin provisioning path: the
service verifies that the caller may grant the requested roles, creates the account with
\texttt{isActive: true}, and attaches the roles.

\reportsideways{figures/fig_3_5}{Create-user-by-admin sequence.}{3-5}

\textbf{Class diagram.} Figure~\ref{fig:3-6} extends the identity core with administrative and
search fields and adds the team graph. \texttt{Team} groups users through \texttt{UserTeam},
whose \texttt{isManager} flag carries management semantics.

\reportfigure[0.90\textwidth]{figures/fig_3_6}{Users \& Teams class diagram.}{3-6}

\subsection{Implementation}

The sprint delivered two backend modules and the guest and administration screens that consume
them.

The \textbf{authentication surface} comprises eight endpoints: \texttt{POST /auths/register},
\texttt{/auths/login}, \texttt{/auths/logout}, the three-step password reset
(\texttt{/auths/request-reset-code}, \texttt{/verify-reset-code}, \texttt{/reset-password}), and
the token endpoints \texttt{/tokens/verify} and \texttt{/tokens/refresh}.

The \textbf{administration surface} comprises fourteen guarded endpoints across \texttt{/users}
and \texttt{/teams}: user provisioning, \texttt{GET /users/me}, the role-scoped user list with
CSV export, self and admin profile updates, password change, soft delete, bulk email, and team
CRUD.

On the \textbf{client}, the guest route group carries the login, registration, and
reset-password screens. The authenticated \texttt{users} section carries the directory table
(with search by name, email, or phone, role filtering, sorting, and pagination), the
create/edit user dialog, and the teams view.

\reportscreenshot{figures/screenshot_P07}{Login page.}{P07}

\reportscreenshot{figures/screenshot_P08}{Registration page.}{P08}

\reportscreenshot{figures/screenshot_P09}{Users list with search and filters.}{P09}

\reportscreenshot{figures/screenshot_P10}{Create/edit user dialog.}{P10}

\reportscreenshot{figures/screenshot_P11}{Teams view.}{P11}

\subsection{Validation and testing}

The testing strategy for this sprint focused on two areas: an \textbf{end-to-end RBAC suite} and
\textbf{acceptance scenario verification}.

The e2e suite (\texttt{test/agile-permissions.e2e-spec.ts}) drives the full permission matrix
over supertest for the CEO, PM, PO, Scrum Master, and engineer roles, verifying that each role
can only reach the routes its permissions allow. This suite validates the core authorization
contract established in this sprint.

The acceptance scenarios below record the Given/When/Then conditions each user story had to
satisfy:

\begin{longtable}{@{}p{1.8cm}p{8.3cm}p{3.8cm}@{}}
\caption{Sprint 1 acceptance scenarios.}\label{tab:s1-acceptance} \\
\tableheaderrow
\thcell{US-ID} & \thcell{Acceptance scenario (Given / When / Then)} & \thcell{Result} \\
\midrule
\endfirsthead
\caption[]{Sprint 1 acceptance scenarios (continued).} \\
\tableheaderrow
\thcell{US-ID} & \thcell{Acceptance scenario (Given / When / Then)} & \thcell{Result} \\
\midrule
\endhead
\bottomrule
\endfoot
\rowcolor{rowalt}
US-S1-01 & Given the backend, When any module is inspected, Then it follows the controller → service → repository → DTO layering over a Prisma multi-file schema with applied migrations. & \pass{} Verified \\
US-S1-02 & Given a valid, active account, When it posts correct email/phone + password, Then a JWT access/refresh pair is returned; wrong password → 401, unknown user → 404. & \pass{} Verified \\
\rowcolor{rowalt}
US-S1-03 & Given a stored refresh token, When \texttt{/tokens/refresh} is called with it, Then a new access token is minted; a never-issued/revoked token is rejected. & \pass{} Verified \\
US-S1-04 & Given a registered email, When the 3-step reset is followed with the emailed 5-digit code inside 15 min, Then the password is changed. & \pass{} Verified \\
\rowcolor{rowalt}
US-S1-05 & Given a protected route with \texttt{@Permissions}, When a caller lacking a matching permission calls it, Then the guard returns 403; a matching role passes. & \pass{} Verified (e2e suite) \\
US-S1-06 & Given an admin holding \texttt{user.create.by.admin}, When they create a user with roles they may grant, Then the account is created active; granting an unmanageable role → 403. & \pass{} Verified \\
\rowcolor{rowalt}
US-S1-07 & Given users exist, When an admin lists with a name/role filter, Then results are accent-insensitively matched and paginated. & \pass{} Verified \\
US-S1-08 & Given a target user, When an admin edits or deletes it, Then edits persist and delete flips \texttt{isActive: false} (soft delete). & \pass{} Verified \\
\rowcolor{rowalt}
US-S1-09 & Given a logged-in user, When they update their profile or change their password (old password re-verified), Then the changes persist. & \pass{} Verified \\
US-S1-10 & Given an admin, When they create a team with members and mark a manager, Then the team and its \texttt{UserTeam} rows (with \texttt{isManager}) are created. & \pass{} Verified \\
\end{longtable}

\subsection{Sprint review}

\textbf{Objectives achieved.} The sprint met its goal: the platform has a working layered
backend, a stateless-JWT authentication core with server-side refresh state, and full user/team
administration governed by a single declarative RBAC catalogue applied uniformly across 139
routes on 18 controllers. The two-tier authorization pattern introduced here --- a coarse
permission guard plus a data-scoped service check --- becomes the platform-wide template every
later module follows.

\textbf{Remaining work.} Security-hardening items identified during the sprint --- shorter token
lifetimes, request throttling, and a stricter authorization check on user deletion --- are logged
and prioritized for future work (Annex~A).

\subsection{Sprint retrospective}

\textbf{Difficulties encountered.} This sprint involved a significant learning curve, as it was
the author's first experience with both NestJS and Prisma. Understanding NestJS's
decorator-based dependency injection, guard composition, and module wiring required upfront
investment before productive development could begin. Similarly, learning Prisma's multi-file
schema, migration workflow, and query API added to the ramp-up time. Designing the RBAC
catalogue --- mapping approximately 120 permissions across 31 role types into a single
declarative constant --- was a non-trivial design effort that required iterating on the
permission granularity to find the right balance between expressiveness and maintainability.

\textbf{Lessons learned.} Investing time in a strict, uniform architecture from the start (the
four-layer controller-service-repository-DTO convention) paid off immediately: once the
pattern was established in this sprint, every subsequent module could be scaffolded quickly
by following the same structure. The table-driven RBAC approach also proved its value early ---
a single guard and one permission constant are far easier to audit than scattered authorization
checks.

\textbf{Improvements planned.} Token TTLs should be shortened before any non-local deployment.
Request throttling should be added to prevent brute-force attacks on the login and
password-reset endpoints. These items are carried forward as prioritized future work.

\subsection{Cumulative class diagram}

Figure~\ref{fig:3-7} is the first cumulative view of the domain model. After Sprint~1 it
contains exactly this sprint's identity core, so every class is new. Later sprints extend this
same diagram, highlighting only their additions, so the reader can watch the model grow.

\reportfigure[\textwidth]{figures/fig_3_7}{Cumulative class diagram after Sprint 1.}{3-7}

\subsection{Sprint conclusion}

Sprint~1 established the foundation the entire platform depends on: a layered backend
architecture, a stateless JWT authentication system, a centralized RBAC model, and user/team
administration. The identity classes and the two-tier authorization pattern introduced here are
the hub that the entire domain model and every subsequent sprint build upon.

% Sprints 2-3 are omitted from the partial (Sprint-1-only) excerpt build.
\ifpartialbuild\else
\clearpage
\section{Sprint 2: Projects \& Membership}

\subsection{Sprint goal}

With identity and access in place, Sprint 2 introduces the unit of work everything else in
the platform hangs off: the \textbf{project}. The goal was the full project lifecycle (create, read, update,
archive, hard-delete), per-project team membership with a single project manager, and an email/token
invitation flow for bringing people who are not yet members onto a project. This is the largest
domain module and the aggregate root of the whole agile domain: epics, sprints, tasks, milestones,
labels and reminders all foreign-key back to a \texttt{Project}, so it had to land before any of the planning or
execution work in the sprints that follow.

\subsection{Sprint backlog}

The user stories for this sprint, repeated from the product backlog in Section~2.2 so the
sprint reads self-contained:

\begin{longtable}{@{}l>{\raggedright\arraybackslash}p{0.50\textwidth}lll@{}}
\caption{Sprint 2 backlog.}\label{tab:s2-backlog-detail} \\
\tableheaderrow
\thcell{US-ID} & \thcell{Story} & \thcell{SP} & \thcell{Est.} & \thcell{Priority} \\
\midrule
\endfirsthead
\caption[]{Sprint 2 backlog (continued).} \\
\tableheaderrow
\thcell{US-ID} & \thcell{Story} & \thcell{SP} & \thcell{Est.} & \thcell{Priority} \\
\midrule
\endhead
\bottomrule
\endfoot
\rowcolor{rowalt}
US-S2-01 & As a manager/executive, I want to create a project under a business unit with a type (AGILE/FREESTYLE), so that work is organized by unit and workflow. & 5 & 2.5d & Must \\
US-S2-02 & As a manager, I want to edit, archive, and delete a project, so that I manage its lifecycle. & 3 & 1.5d & Must \\
\rowcolor{rowalt}
US-S2-03 & As a manager, I want to add existing users as project members and mark a manager, so that the team is defined. & 3 & 1.5d & Must \\
US-S2-04 & As a manager, I want to invite people by email with a single-use token, so that non-members can be brought in. & 5 & 3d & Should \\
\rowcolor{rowalt}
US-S2-05 & As an invitee, I want to accept an emailed invitation, so that I become a member. & 3 & 1.5d & Should \\
US-S2-06 & As a manager, I want to configure the project's kanban settings (WIP limits), so that the board fits our process. & 2 & 1d & Could \\
\rowcolor{rowalt}
US-S2-07 & As a member, I want to see only the projects I belong to (executives scoped to their business unit), so that access is isolated. & 5 & 2d & Must \\
\midrule
 & \textbf{Sprint 2 subtotal} & \textbf{26} & \textbf{2w} & \\
\end{longtable}

\subsection{Requirements analysis}

The sprint delivers one module. The features, actors, and use cases are collected in
Figure~\ref{fig:3-8}.

\reportfigure[0.70\textwidth]{figures/fig_3_8}{Sprint 2 use-case diagram.}{3-8}

Three actor archetypes appear. An \textbf{Executive} (CEO, CTO, or CMO) creates projects and
manages them globally or within their business unit: the \texttt{CEO} sees every project, while the
\texttt{CTO} and \texttt{CMO} are scoped to \texttt{TawerDev} and \texttt{TawerCreative} respectively. A
\textbf{Project Manager} --- the single \texttt{ProjectMember} marked \texttt{isManager} --- manages the team,
invites by email, and configures the board, but only for projects they manage. A \textbf{Member}
reads the projects they belong to.

\subsubsection{Projects, Membership \& Invitations}

\textbf{Functional description.} A project is created by an executive, belongs to exactly one
business unit (\texttt{TawerDev} or \texttt{TawerCreative}), and carries a \texttt{projectType} of
\texttt{AGILE} or \texttt{FREESTYLE} that decides which planning features apply in later sprints. It
gathers a team through \texttt{ProjectMember} rows, of which exactly one is flagged \texttt{isManager};
that single-manager invariant is enforced on create and protected on every mutation, so demoting
or removing the last manager is rejected. Translatable fields (\texttt{name}, \texttt{description},
\texttt{details}) live in a child \texttt{ProjectContent} table keyed by language, mirroring the
app-wide i18n content pattern. WIP limits per kanban column are stored as a JSON map on the
project itself (\texttt{kanbanSettings}), validated against the project's actual status names.
Two onboarding paths exist: executives can attach an existing user to a project directly, while
managers invite by email through a tokenized \texttt{ProjectInvitation} with a default seven-day
expiry.

\textbf{Actors.} An Executive initiates project creation and can add members directly. A Project
Manager manages the team and invites by email. A Member reads the projects they belong to. An
Invitee (not yet a member) accepts a tokenized invitation to join.

\textbf{Business rules.}
\begin{itemize}
  \item Authorization reuses the \textbf{two-tier pattern} from Sprint~1: the coarse
    \texttt{HasPermissionGuard} gates the route, then a fine-grained check runs inside the service.
    The read isolation is pushed into the Prisma \texttt{where} clause and runs atomically with the
    fetch, so a caller who lacks access simply gets zero rows or a \texttt{403}.
  \item A non-executive caller is pinned to projects they belong to
    (\texttt{members.some(\{ userId \})}); a unit executive is pinned to their \texttt{businessUnit};
    the CEO is global; management operations additionally require \texttt{isManager: true} or
    executive status.
  \item The \textbf{single-manager invariant} is enforced in three places: on create (exactly one
    manager required), on demotion (the last manager cannot be demoted), and on removal (the last
    manager cannot be removed by a non-executive).
  \item The module uses the Prisma query-builder exclusively, with no raw SQL, so there is no
    string-interpolation injection surface.
  \item The email invitation path verifies eligibility (not already a member or invited), upserts
    a \texttt{ProjectInvitation} carrying a \texttt{randomUUID} token, and on acceptance re-verifies
    that the accepting user's email matches the invited email, so one user cannot claim another's
    invitation.
\end{itemize}

\textbf{Create-project sequence.} Figure~\ref{fig:3-9} shows the create path: the service
re-checks executive status and business-unit access (defense in depth), normalizes the input and
defaults the manager to the creator when none is given, validates the single-manager invariant,
then writes the project with its nested members and contents in one repository call; a duplicate
name surfaces as a Prisma \texttt{P2002} that becomes a \texttt{409}.

% Figure 3.9 — render diagrams-src/fig_3_9.puml -> figures/fig_3_9.(png|pdf)
\reportfigure{figures/fig_3_9}{Create-project sequence diagram.}{3-9}

\textbf{Add-member / invite sequence.} Figure~\ref{fig:3-10} shows the single ``smart'' endpoint
(\texttt{POST /:projectId/members}) that branches on whether the caller supplied a \texttt{userId} or
an \texttt{email}. The \texttt{userId} branch is executive-only and adds the member directly. The
\texttt{email} branch is open to executives and the project's manager: it upserts a tokenized
invitation, sends a branded email with a join link, and raises an in-app notification if the
invitee already has an account.

% Figure 3.10 — render diagrams-src/fig_3_10.puml -> figures/fig_3_10.(png|pdf)
\reportsideways{figures/fig_3_10}{Add-member / invite sequence diagram.}{3-10}

\textbf{Class diagram.} Figure~\ref{fig:3-11} adds the project family to the model. \texttt{Project}
is the aggregate root; its translatable fields are split into \texttt{ProjectContent};
\texttt{ProjectMember} is the join to \texttt{User} that also carries the \texttt{isManager} flag and a
per-member \texttt{hourlyRate}; and \texttt{ProjectInvitation} holds the email-bound token lifecycle.
\texttt{User} is shown again --- not new --- because it now gains three roles against the project
family: creator, member, and inviter.

% Figure 3.11 — render diagrams-src/fig_3_11.puml -> figures/fig_3_11.(png|pdf)
{\renewcommand{\figmaxheight}{0.62\textheight}%
\reportfigure[\textwidth]{figures/fig_3_11}{Projects \& Membership class diagram.}{3-11}}

\subsection{Implementation}

The sprint delivered one backend module and the project management screens that consume it.

The \textbf{project surface} comprises seventeen endpoints on \texttt{ProjectsController} covering
the lifecycle (register, list, get, get capacity, update, delete, archive, restore), membership
(smart add-member, update-member, remove-member), the invitation flow (create, revoke, resend,
accept), and the kanban settings (get, patch).

On the \textbf{client}, the projects section is a status-tabbed list with search, a filter panel,
and drag-reorder. Opening a project reveals a tabbed detail whose \textbf{Members} tab lists members
and pending invitations with direct-add, invite-by-email, role-toggle, remove, and resend/revoke
actions.

\reportscreenshot{figures/screenshot_P12}{Projects list view.}{P12}

\reportscreenshot{figures/screenshot_P13}{Create-project sheet.}{P13}

\reportscreenshot{figures/screenshot_P14}{Members tab.}{P14}

\reportscreenshot{figures/screenshot_P15}{Invite-by-email dialog.}{P15}

\subsection{Validation and testing}

The testing strategy for this sprint focused on two areas: a \textbf{behavioural unit-test suite}
and \textbf{acceptance scenario verification}.

The unit-test suite (\texttt{projects.service.spec.ts}, 854 lines) exercises the full project
lifecycle with mocked repositories: create (forbidden, success, duplicate-name \texttt{P2002}),
the executive-vs-non-executive list scoping, get-by-id (member, executive, not-found), capacity,
update (business-unit-immutable, success, conflict), delete (including the CTO business-unit
filter), archive/restore, and add/update/remove-member. A companion
\texttt{projects.controller.spec.ts} (215 lines) validates the controller wiring.

The acceptance scenarios below record the Given/When/Then conditions each user story had to
satisfy:

\begin{longtable}{@{}p{1.8cm}p{8.3cm}p{3.8cm}@{}}
\caption{Sprint 2 acceptance scenarios.}\label{tab:s2-acceptance} \\
\tableheaderrow
\thcell{US-ID} & \thcell{Acceptance scenario (Given / When / Then)} & \thcell{Result} \\
\midrule
\endfirsthead
\caption[]{Sprint 2 acceptance scenarios (continued).} \\
\tableheaderrow
\thcell{US-ID} & \thcell{Acceptance scenario (Given / When / Then)} & \thcell{Result} \\
\midrule
\endhead
\bottomrule
\endfoot
\rowcolor{rowalt}
US-S2-01 & Given an executive, When they POST a project with a business unit and type, Then a Project with nested contents and one manager member is created; a non-executive or wrong-unit caller \textrightarrow{} 403; a duplicate name \textrightarrow{} 409. & \pass{} Verified \\
US-S2-02 & Given a manager/executive, When they PATCH, archive, or DELETE a project, Then fields update in a transaction, archive flips \texttt{isArchived}, and delete cascades the whole project. & \pass{} Verified \\
\rowcolor{rowalt}
US-S2-03 & Given an executive, When they add an existing user by \texttt{userId}, Then a \texttt{ProjectMember} is created unless the user is already a member (409). & \pass{} Verified \\
US-S2-04 & Given an executive or the project manager, When they invite an email, Then a single-use tokenized invitation (7-day expiry) is upserted, emailed, and notified if the invitee has an account. & \pass{} Verified \\
\rowcolor{rowalt}
US-S2-05 & Given a valid, unexpired invitation, When the authenticated invitee (email matching) accepts, Then a membership is created and the invite marked ACCEPTED; a mismatched email or expired token is rejected. & \pass{} Verified \\
US-S2-06 & Given a manager, When they PATCH kanban settings, Then WIP-limit keys are validated against the project's status names and stored as JSON. & \pass{} Verified \\
\rowcolor{rowalt}
US-S2-07 & Given a non-executive, When they list projects, Then only projects they are a member of are returned; a CTO/CMO sees only their business unit; a CEO sees all. & \pass{} Verified \\
\end{longtable}

\subsection{Sprint review}

\textbf{Objectives achieved.} The sprint met its goal: the platform has a project aggregate root
with full lifecycle management, business-unit scoping enforced inside the query rather than bolted
on after it, a single-manager invariant guarded on every mutation, and an email-based invitation
flow with identity-verified token acceptance. It is also the first module with a dedicated
behavioural test suite --- an 854-line \texttt{projects.service.spec.ts} exercising the
create/list/update/delete/membership paths --- which raised confidence in exactly those paths.

\textbf{Remaining work.} Hardening items identified during the sprint --- server-side
archived-project filtering on the list endpoint, a project-scoped name uniqueness constraint, and
an independent per-member capacity input --- are logged and prioritized for future work (Annex~A).

\subsection{Sprint retrospective}

\textbf{Difficulties encountered.} The main challenge was designing the project-scoped
authorization model. Unlike Sprint~1's global RBAC, this sprint required data-scoped isolation
where the permission predicate and the database query had to be a single atomic statement. Getting
the business-unit scoping right --- CEO global, CTO/CMO unit-pinned, manager project-pinned ---
required careful layering of the Prisma \texttt{where} clauses. The single-manager invariant also
demanded attention: enforcing it across create, demotion, and removal without leaving an
inconsistent window required validating at three separate points in the service.

\textbf{Lessons learned.} Pushing the access check into the query itself (rather than fetching
first and filtering after) proved to be the strongest pattern in the module: it eliminates the
window between a check passing and data escaping its scope. The i18n content-table pattern
(\texttt{ProjectContent} keyed by language) also demonstrated its reusability --- it was adopted
directly in later sprints for sprints and tasks.

\textbf{Improvements planned.} Archived-project filtering should move to the server side so the
list endpoint excludes them by default. The project-name uniqueness constraint should be scoped
per project rather than enforced globally.

\subsection{Cumulative class diagram}

Figure~\ref{fig:3-12} focuses on the entities introduced in this sprint. The \texttt{User} anchor
from Sprint~1 is shown in plain style so the attachment point is clear; this sprint's additions ---
\texttt{Project}, \texttt{ProjectContent}, \texttt{ProjectMember}, \texttt{ProjectInvitation} and
their enums --- are highlighted. The link back to \texttt{User} (creator, member, inviter) is the
seam that ties the new project aggregate to the identity core established previously.

% Figure 3.12 — render diagrams-src/fig_3_12.puml -> figures/fig_3_12.(png|pdf)
{\renewcommand{\figmaxheight}{0.82\textheight}%
\reportfigure[\textwidth]{figures/fig_3_12}{Cumulative class diagram after Sprint 2.}{3-12}}

\subsection{Sprint conclusion}

Sprint~2 delivered the aggregate root of the entire agile domain: the project container with
lifecycle management, business-unit scoping, team membership with a guarded single-manager
invariant, and a tokenized email invitation flow. The query-scoped authorization pattern
introduced here --- merging the access check into the database query itself --- and the i18n
content-table pattern are reused directly by every subsequent sprint.

\clearpage
\section{Sprint 3: Agile Backlog \& Tasks}

\subsection{Sprint goal}

With people, access, and projects in place, Sprint~3 adds the work itself. The sprint delivers
two inseparable modules: the agile planning layer (epics, sprints with a governed lifecycle,
milestones, and burndown/velocity/Gantt analytics) and, on top of it, the task engine --- a
per-project data-driven Kanban with custom columns, WIP limits, dependency-blocking, story
points, time logging, labels, and threaded comments. They land together because a sprint is only
meaningful once tasks can be assigned to it, and the burndown chart reads directly from each
task's \texttt{storyPoints} and \texttt{completedAt}. At \textbf{64 story points} over five weeks
this is the heaviest load in the plan and the richest domain code in the codebase.

\subsection{Sprint backlog}

The user stories for this sprint, repeated from the product backlog in Section~2.2 so the
sprint reads self-contained:

\begin{longtable}{@{}l>{\raggedright\arraybackslash}p{0.50\textwidth}lll@{}}
\caption{Sprint 3 backlog.}\label{tab:s3-backlog-detail} \\
\tableheaderrow
\thcell{US-ID} & \thcell{Story} & \thcell{SP} & \thcell{Est.} & \thcell{Priority} \\
\midrule
\endfirsthead
\caption[]{Sprint 3 backlog (continued).} \\
\tableheaderrow
\thcell{US-ID} & \thcell{Story} & \thcell{SP} & \thcell{Est.} & \thcell{Priority} \\
\midrule
\endhead
\bottomrule
\endfoot
\rowcolor{rowalt}
US-S3-01 & As a PM/PO, I want to create epics to group large features, so that the backlog is structured. & 3 & 1.5d & Should \\
US-S3-02 & As a Scrum Master, I want to create sprints with dates and a capacity, so that work is time-boxed. & 5 & 3d & Must \\
\rowcolor{rowalt}
US-S3-03 & As a Scrum Master, I want to run the sprint lifecycle (start/stop/complete, one running at a time), so that iterations are governed. & 8 & 4d & Must \\
US-S3-04 & As a PM, I want milestones with target dates (also for FREESTYLE projects), so that deadlines are tracked. & 3 & 1.5d & Should \\
\rowcolor{rowalt}
US-S3-05 & As a PM, I want burndown, velocity, and Gantt analytics, so that I can monitor progress. & 8 & 3.5d & Should \\
US-S3-06 & As a member, I want to create tasks with type, priority, assignee, and estimates, so that work is captured. & 5 & 2.5d & Must \\
\rowcolor{rowalt}
US-S3-07 & As a manager, I want a per-project, data-driven kanban with custom columns and WIP limits, so that the board mirrors our workflow. & 8 & 4d & Must \\
US-S3-08 & As an assignee, I want to move a task across columns with transition, dependency, and WIP validation, so that status stays valid. & 8 & 3.5d & Must \\
\rowcolor{rowalt}
US-S3-09 & As a member, I want to declare task dependencies (blocking / blocked-by), so that ordering is explicit. & 5 & 2d & Should \\
US-S3-10 & As a member, I want to log time against a task, so that effort is recorded. & 3 & 1.5d & Should \\
\rowcolor{rowalt}
US-S3-11 & As a member, I want threaded comments with @mentions and likes on a task, so that we collaborate in context. & 5 & 2.5d & Should \\
US-S3-12 & As a member, I want labels and to assign a task to a sprint/epic/milestone, so that work is categorized. & 3 & 1.5d & Should \\
\midrule
 & \textbf{Sprint 3 subtotal} & \textbf{64} & \textbf{5w} & \\
\end{longtable}

\subsection{Requirements analysis}

The sprint delivers two modules. The features, actors, and use cases are collected in
Figure~\ref{fig:3-13}.

{\renewcommand{\figmaxheight}{0.75\textheight}%
\reportfigure[\textwidth]{figures/fig_3_13}{Sprint 3 use-case diagram.}{3-13}}

Four Scrum-shaped roles appear, each mapped to a capability helper the service enforces. A
\textbf{Project Manager} and \textbf{Product Owner} shape structure --- epics, milestones, task
creation, dependencies, labels, and board columns. A \textbf{Scrum Master} governs the sprint
lifecycle and, with the PM/PO, grooms the backlog and assigns items to sprints. Any
\textbf{Assignee/Developer} who is a project member performs the day-to-day work --- moving their
own cards, logging time, and commenting.

\subsubsection{Module A --- Agile Backlog (Epics, Sprints \& Milestones)}

\textbf{Functional description.} The agile layer adds three planning artefacts, each scoped to a
project. An \textbf{epic} is a large feature that groups tasks under a named, colour-coded,
optionally date-bounded umbrella so progress can be rolled up. A \textbf{sprint} is a time-boxed
iteration with planned and estimated date windows, a story-point \texttt{capacity}, multilingual
content, file attachments, and a genuine lifecycle. A \textbf{milestone} is a target date for a
set of tasks. Epics and sprints are hard-gated to \texttt{AGILE} projects by an
\texttt{AgileOnlyGuard}, while milestones carry no such guard because a deadline is useful to a
\texttt{FREESTYLE} project too. All three follow the project-wide four-layer pattern
(controller → service → repository → DTO) with four thin repositories apiece, and all three
enqueue an AI embedding job on every write via the \texttt{IndexOutboxService}, so agile artefacts
are first-class documents for the copilot built in Sprint~6.

\textbf{Actors.} A \textbf{Scrum Master} creates and manages sprints and controls the sprint
lifecycle (start, stop, complete, restart). A \textbf{Product Owner} creates and manages epics and
milestones. A \textbf{Project Manager} accesses the burndown, velocity, and Gantt analytics.
Authorization reuses the two-tier model from earlier sprints --- a coarse
\texttt{HasPermissionGuard} at the route, then a project-scoped capability check inside the
service.

\textbf{Business rules.}
\begin{itemize}
  \item Epics and sprints are \textbf{gated to \texttt{AGILE} projects} by an
    \texttt{AgileOnlyGuard}; milestones are available to both project types.
  \item A sprint follows a \textbf{state machine}: created \texttt{Pending}, started to
    \texttt{Running} (guarded by the rule that no other sprint in the same project may already be
    running), and from \texttt{Running} it may be stopped or completed. Both \texttt{Stopped} and
    \texttt{Completed} can be restarted back to \texttt{Running}.
  \item Moving to \texttt{Running} or \texttt{Completed} fans out a notification to every project
    member; moving to \texttt{Stopped} or \texttt{Completed} cancels the sprint's still-pending
    reminders.
  \item Sprint date windows are validated: end after start, estimated-end after estimated-start,
    and the whole sprint window must sit inside the project's own date window.
  \item Sprint content supports multilingual entries (the same i18n content-table pattern used for
    projects) and file attachments through a child \texttt{SprintAttachment} collection.
  \item Names are accent-stripped into an \texttt{unaccentedName} field for search, mirroring the
    pattern established for users and projects.
\end{itemize}

\textbf{Create-sprint sequence.} Figure~\ref{fig:3-14} shows the module's most involved write
path. After the RBAC and AGILE guards pass, the service runs its own \texttt{canManageProject}
check, requires at least one content entry, stamps the accent-stripped name, and validates every
date window. It then persists the sprint in the \texttt{Pending} state with its nested content and
attachment rows and creates up to two default reminders (one day before start, two days before
end). Finally, it enqueues the embedding upsert so the sprint is immediately searchable.

% Figure 3.14 — render diagrams-src/fig_3_14.puml -> figures/fig_3_14.(png|pdf)
\reportsideways{figures/fig_3_14}{Create-sprint sequence diagram.}{3-14}

\textbf{Sprint analytics sequence.} The planning artefacts pay off as analytics, computed
server-side from real task data. Burndown (Figure~\ref{fig:3-15}) is the clearest example: for a
chosen sprint the service loads every task's \texttt{storyPoints}, \texttt{status},
\texttt{completedAt} and \texttt{createdAt}, sums the committed points, then walks day-by-day
across the sprint window computing an \emph{ideal} line (total minus an even daily burn) against
the \emph{actual} remaining (total minus the points of tasks marked \texttt{DONE} on or before that
day). Velocity aggregates completed-sprint points across the project, and the Gantt view fans
four parallel queries --- milestones, epics, sprints, tasks --- into one timeline payload.

% Figure 3.15 — render diagrams-src/fig_3_15.puml -> figures/fig_3_15.(png|pdf)
\reportsideways{figures/fig_3_15}{Sprint burndown sequence diagram.}{3-15}

\textbf{Sprint lifecycle state machine.} The richest piece of business logic in this module is
the sprint state machine (Figure~\ref{fig:3-16}). A sprint is always created \texttt{Pending}.
Starting it is guarded by the rule that no other sprint in the same project may already be
running. From \texttt{Running} a sprint may be stopped or completed, and both \texttt{Stopped} and
\texttt{Completed} can be restarted back to \texttt{Running}. Two side effects hang off the
transitions: moving to \texttt{Running} or \texttt{Completed} fans out a notification to every
other project member, and moving to \texttt{Stopped} or \texttt{Completed} cancels the sprint's
still-pending reminders. The state machine is a shared contract --- the backend validator and the
frontend sprint card's action buttons (Start / Stop / Complete / Restart) mirror each other
exactly.

% Figure 3.16 — render diagrams-src/fig_3_16.puml -> figures/fig_3_16.(png|pdf)
\reportfigure[0.65\textwidth]{figures/fig_3_16}{Sprint lifecycle state machine.}{3-16}

\textbf{Class diagram.} Figure~\ref{fig:3-17} introduces the five planning entities and the
\texttt{SprintStatus} enum. \texttt{Sprint} splits its translatable name and description into a
child \texttt{SprintContent} (the same i18n content-table pattern used for projects) and owns a
\texttt{SprintAttachment} collection for files. \texttt{Epic} and \texttt{Milestone} sit directly
under the project. All three carry the optional foreign keys that \texttt{Task} (defined in
Module~B) will point back through --- \texttt{epicId}, \texttt{sprintId}, \texttt{milestoneId} ---
which is why the diagram shows the relations to \texttt{Task} even though the class itself belongs to
the next module.

% Figure 3.17 — render diagrams-src/fig_3_17.puml -> figures/fig_3_17.(png|pdf)
\reportfigure[0.90\textwidth]{figures/fig_3_17}{Agile Backlog class diagram.}{3-17}

\subsubsection{Module B --- Tasks \& Data-Driven Kanban}

\textbf{Functional description.} The task engine is the core unit of work and the richest domain
module in the codebase (\texttt{tasks.service.ts} alone is 2\,735 lines). A task carries the
expected fields --- type, priority, assignee, reporter, story points, estimated and actual hours,
due date --- plus the machinery for a real board: a free-string \texttt{status}, a
\texttt{displayOrder} for manual ordering, one level of subtasks, dependencies, labels, time
entries, and threaded comments with likes and mentions. The design adapts to the project type: an
\texttt{AGILE} project gets the full six-column workflow (\texttt{BACKLOG → TODO → IN\_PROGRESS →
IN\_REVIEW → TESTING → DONE}) with sprints, epics and story points, while a \texttt{FREESTYLE}
project gets a lean three-column board (\texttt{TODO / IN\_PROGRESS / DONE}) and a progress
percentage instead; the service enforces which fields are allowed per type.

The board itself is \textbf{data-driven}. Rather than hard-code columns, each project owns a set
of \texttt{ProjectTaskStatus} rows --- the source of truth for its columns, their order, colours,
and the legal transitions out of each. A task's \texttt{status} is a plain string that holds either
a system enum name or a custom column name, which lets a team invent columns without a schema
change; the link between \texttt{Task.status} and \texttt{ProjectTaskStatus.name} is
application-enforced since a free string cannot carry a foreign key. Statuses are seeded lazily on
first read of the board rather than at project creation, and transition validation prefers the DB
rows, falling back to hard-coded maps only when a project has no statuses yet. WIP limits are
stored in the project's \texttt{kanbanSettings} JSON field and validated against the project's
actual status names on every column move.

\textbf{Actors.} A \textbf{Project Manager / Product Owner} creates tasks, manages structure
(dependencies, labels, board columns), and grooms the backlog. A \textbf{Scrum Master} assigns
tasks to sprints and manages sprint-level backlog operations. An \textbf{Assignee / Developer}
moves their own cards across columns, logs time, and comments. Authorization is graduated into
seven capability helpers --- from \texttt{canAccessProject} (any member may read, comment, and log
time) up through \texttt{canCreateTaskForProject}, \texttt{canManageTaskStructure},
\texttt{canManageBacklog}, \texttt{canManageSprintAssignment}, and
\texttt{canAdvanceTaskWorkflow}, the last of which lets the assignee move their own card even
without a management role.

\textbf{Business rules.}
\begin{itemize}
  \item An \texttt{AGILE} project gets a \textbf{six-column workflow}
    (\texttt{BACKLOG → TODO → IN\_PROGRESS → IN\_REVIEW → TESTING → DONE}); a \texttt{FREESTYLE}
    project gets a \textbf{three-column board} (\texttt{TODO / IN\_PROGRESS / DONE}). The service
    rejects fields that do not apply to the project type (e.g.\ a sprint or epic reference on a
    \texttt{FREESTYLE} task).
  \item Board columns are \textbf{data-driven}: \texttt{ProjectTaskStatus} rows define the columns,
    their order, and the legal transitions out of each. A team can add custom columns without a
    schema change.
  \item A column move is validated against \textbf{three gates} in order: the transition must be
    legal for the project's status rows, the task must not be blocked by any dependency that is not
    yet \texttt{DONE}, and the target column must be under its WIP limit. Each gate maps to a
    distinct error (\texttt{TASK\_BLOCKED}, \texttt{WIP\_LIMIT\_REACHED}).
  \item Moving a task to \texttt{DONE} stamps \texttt{completedAt}; the assignee is notified of
    every status change via a fire-and-forget notification.
  \item A \textbf{circular-dependency guard} rejects cycles when adding a dependency: a recursive
    traversal of the blocking graph detects self-references, bidirectional edges, and transitive
    cycles.
  \item Each task gets a generated \texttt{TASK-<n>} key scoped to the project.
  \item Task comments are created in a transaction with their \texttt{@mentions}; mentioned users
    and the assignee are notified.
  \item Every task and comment write enqueues an AI embedding upsert via the
    \texttt{IndexOutboxService}, so the entire task corpus is searchable ahead of Sprint~6.
\end{itemize}

\textbf{Create-task sequence.} Figure~\ref{fig:3-18} shows the create path, which validates on
three axes before it writes: the caller's capability, the project-type field rules (rejecting a
sprint, epic, or story-point value on a \texttt{FREESTYLE} task), and referential integrity (any
sprint, epic, or milestone named must belong to this project). The repository then generates a
\texttt{TASK-<n>} key and inserts the task with its attachments. Afterwards the service optionally
spawns due-date reminders, fires a fire-and-forget assignment notification if the assignee is not
the author, and enqueues the AI index upsert.

% Figure 3.18 — render diagrams-src/fig_3_18.puml -> figures/fig_3_18.(png|pdf)
\reportsideways{figures/fig_3_18}{Create-task sequence diagram.}{3-18}

\textbf{Move-task sequence.} The Kanban move (Figure~\ref{fig:3-19}) is where the workflow engine
enforces its constraints. A drag-drop issues a single PATCH; the service loads the task, checks
\texttt{canAdvanceTaskWorkflow} (assignee or a manager/PO/Scrum-Master/executive), then runs three
gates in order: the transition must be legal for the project's status rows, the task must not be
blocked by any dependency that is not yet \texttt{DONE}, and the target column must be under its
WIP limit. Only if all three pass does it update the status --- stamping \texttt{completedAt} when
the task lands on \texttt{DONE} --- and fire the assignee notification.

% Figure 3.19 — render diagrams-src/fig_3_19.puml -> figures/fig_3_19.(png|pdf)
\reportsideways{figures/fig_3_19}{Move-task Kanban sequence diagram.}{3-19}

\textbf{Task status transitions.} The status transitions (Figure~\ref{fig:3-20}) show the two
seeded system boards. \texttt{FREESTYLE} keeps a simple three-column loop; \texttt{AGILE} runs the
full six-column pipeline where each column allows moving forward or one step back (a card in
\texttt{TESTING} can drop to \texttt{IN\_REVIEW} or advance to \texttt{DONE}, and \texttt{DONE} can
reopen to \texttt{TESTING}). These are the default transitions seeded per project; a team can edit
its \texttt{ProjectTaskStatus} rows to add custom columns, and a move into a custom status is
permitted from anywhere, keeping the board flexible while still guarding moves between the known
system columns.

% Figure 3.20 — render diagrams-src/fig_3_20.puml -> figures/fig_3_20.(png|pdf)
{\renewcommand{\figmaxheight}{0.75\textheight}%
\reportfigure[\textwidth]{figures/fig_3_20}{Task status transitions for FREESTYLE and AGILE boards.}{3-20}}

\textbf{Class diagram.} Figure~\ref{fig:3-21} is the largest single-module class diagram in the
report. \texttt{Task} is the hub; \texttt{ProjectTaskStatus} holds the board columns and their
\texttt{allowedTransitions}; \texttt{TaskComment} carries \texttt{TaskCommentLike} and
\texttt{TaskCommentMention} children; \texttt{TaskLabel} reaches \texttt{Task} through the
\texttt{TaskLabelAssignment} junction; \texttt{TaskDependency} is the self-referential blocking
relation (a task both blocks and is blocked by others); and \texttt{TaskTimeEntry} records logged
hours, optionally linked to a work session (introduced in Sprint~4). Two enums, \texttt{TaskType}
and \texttt{TaskPriority}, constrain the classification fields; the \texttt{status} field is
deliberately left as a free string because the board is data-driven.

% Figure 3.21 — render diagrams-src/fig_3_21.puml -> figures/fig_3_21.(png|pdf)
\reportsideways{figures/fig_3_21}{Tasks \& Kanban class diagram.}{3-21}

\subsection{Implementation}

Between the two modules the sprint delivers a large API surface --- nineteen agile endpoints
across the epics, sprints and milestones controllers (including the burndown, velocity and Gantt
analytics routes), and thirty-five task endpoints covering task CRUD, the dynamic status board,
Kanban move, backlog reorder and move-to-sprint, dependencies, time entries, comments, and
labels. On the client all of this lives inside the \textbf{project-detail} page rather than in
standalone routes: the project opens onto Tasks, Kanban, Backlog, Sprints and Milestones tabs, a
card click opens a detail sheet composed of comment, attachment, dependency, label, subtask and
time-entry sections, and the analytics tab renders the burndown and velocity charts and the Gantt.

\reportscreenshot{figures/screenshot_P16}{Backlog view.}{P16}
\reportscreenshot{figures/screenshot_P17}{Kanban board.}{P17}
\reportscreenshot{figures/screenshot_P18}{Task detail sheet.}{P18}
\reportscreenshot{figures/screenshot_P19}{Sprint board.}{P19}
\reportscreenshot{figures/screenshot_P20}{Burndown chart.}{P20}
\reportscreenshot{figures/screenshot_P21}{Velocity chart.}{P21}
\reportscreenshot{figures/screenshot_P22}{Gantt chart.}{P22}

\subsection{Validation and testing}

The testing strategy for this sprint focused on two areas: \textbf{property-based frontend
testing} and \textbf{acceptance scenario verification}.

Eleven property-based test suites (\texttt{fast-check} + \texttt{vitest}) pin client-side
invariants across the sprint's transform and parsing layers: Kanban no-task-loss and swimlane
membership, backlog casting, bulk-status transforms, time-entry and dependency normalization,
comment casting, analytics transforms, epic and milestone casting, label casting, and task-status
validation. Property-based testing generates hundreds of random inputs per run and asserts an
invariant holds across all of them, which is a stronger guarantee than a handful of example
cases.

The acceptance scenarios below record the Given/When/Then conditions each user story had to
satisfy:

\begin{longtable}{@{}p{1.8cm}p{8.3cm}p{3.8cm}@{}}
\caption{Sprint 3 acceptance scenarios.}\label{tab:s3-acceptance} \\
\tableheaderrow
\thcell{US-ID} & \thcell{Acceptance scenario (Given / When / Then)} & \thcell{Result} \\
\midrule
\endfirsthead
\caption[]{Sprint 3 acceptance scenarios (continued).} \\
\tableheaderrow
\thcell{US-ID} & \thcell{Acceptance scenario (Given / When / Then)} & \thcell{Result} \\
\midrule
\endhead
\bottomrule
\endfoot
\rowcolor{rowalt}
US-S3-01 & Given a PM/PO on an AGILE project, When they create an epic, Then it is stored under the project (unique name per project) and its progress rolls up from DONE task counts; a linked-task epic cannot be deleted. & \pass{} Verified \\
US-S3-02 & Given a Scrum Master, When they create a sprint with all four dates and a capacity, Then a \texttt{Pending} sprint is stored only if every window sits inside the project window; default reminders are created. & \pass{} Verified \\
\rowcolor{rowalt}
US-S3-03 & Given a \texttt{Pending}/\texttt{Running} sprint, When the SM starts/stops/completes it, Then only legal transitions apply, at most one sprint runs per project, members are notified, and pending reminders cancel on stop/complete. & \pass{} Verified \\
US-S3-04 & Given a project (AGILE or FREESTYLE), When a PM creates a milestone with a due date, Then it is stored and can be completed by stamping \texttt{completedAt}. & \pass{} Verified \\
\rowcolor{rowalt}
US-S3-05 & Given a sprint/project with task points, When a PM opens analytics, Then burndown shows ideal-vs-actual, velocity sums completed-sprint points, and the Gantt fans in milestones/epics/sprints/tasks. & \pass{} Verified \\
US-S3-06 & Given a member with create rights, When they POST a task with type/priority/assignee/estimate, Then it is created with a generated key after project-type field and FK validation; assignee is notified; the task is indexed. & \pass{} Verified \\
\rowcolor{rowalt}
US-S3-07 & Given a manager, When they define board columns and WIP limits, Then \texttt{ProjectTaskStatus} rows drive the board and WIP keys are validated against status names. & \pass{} Verified \\
US-S3-08 & Given an assignee, When they drag a card to a new column, Then the move is rejected unless the transition is legal, no blocking dependency is unfinished, and the column is under its WIP limit; \texttt{completedAt} stamps on DONE. & \pass{} Verified \\
\rowcolor{rowalt}
US-S3-09 & Given a member, When they add a dependency, Then a circular-dependency guard rejects cycles; a task cannot leave its column while blocked. & \pass{} Verified \\
US-S3-10 & Given a member, When they log time on a task, Then a \texttt{TaskTimeEntry} is created and \texttt{actualHours} is re-aggregated from the sum. & \pass{} Verified \\
\rowcolor{rowalt}
US-S3-11 & Given a member, When they comment with @mentions, Then the comment and its mentions are created in a transaction, mentioned users and the assignee are notified, and the comment is indexed; likes toggle one-per-user. & \pass{} Verified \\
US-S3-12 & Given a member, When they create/assign labels or move a task to a sprint/epic/milestone, Then project-scoped labels attach via the junction and the FK is set. & \pass{} Verified \\
\end{longtable}

\subsection{Sprint review}

\textbf{Objectives achieved.} The sprint delivered the substance of the platform: a governed
sprint lifecycle whose state machine is a shared frontend/backend contract, server-side
burndown/velocity/Gantt analytics computed from real task points, and a data-driven Kanban that
is a genuine workflow engine --- custom columns, per-column WIP limits, dependency-blocking, and
a circular-dependency guard, all applied under a graduated, project-scoped capability model that
even lets an assignee advance their own card. Every mutation feeds the AI index outbox, so the
whole agile and task corpus is already searchable ahead of Sprint~6.

\textbf{Remaining work.} Hardening items identified during the sprint --- backend service tests
for the tasks module, transaction-wrapped agile writes, and a shared
\texttt{ProjectAccessService} --- are logged and prioritized for future work (Annex~A).

\subsection{Sprint retrospective}

\textbf{Difficulties encountered.} This was the heaviest sprint in the plan (64 story points over
five weeks) and the most demanding in domain complexity. The data-driven Kanban engine required
designing a system where board columns, their transitions, and WIP limits are stored as data
rather than code, which introduced application-level invariants that could not be enforced by the
database schema alone. The graduated capability model --- seven distinct helpers from read-only
access up to workflow advancement, with an assignee-specific override --- required careful layering
to avoid either over-permitting or locking out legitimate moves. The sprint lifecycle state
machine, with its notification side effects and reminder cancellation on certain transitions,
added further coordination between subsystems.

\textbf{Lessons learned.} Property-based testing proved its value for the transform and parsing
layers: generating hundreds of random inputs per run caught edge cases that hand-written examples
would have missed, particularly in the Kanban no-task-loss invariant and dependency normalization.
The fire-and-forget notification pattern (unawaited calls to the notification service) kept the
critical write path fast without sacrificing user feedback. The AI embedding outbox --- enqueueing
an index upsert on every agile and task write --- was straightforward to wire in and ensures the
entire corpus is searchable by the time the copilot lands in Sprint~6.

\textbf{Improvements planned.} Backend service tests for the tasks module's transition
validation, capability matrix, and circular-dependency guard. A shared
\texttt{ProjectAccessService} to replace the duplicated executive-scope helpers across the agile
services. These items are carried forward as prioritized future work (Annex~A).

\subsection{Cumulative class diagram}

Figure~\ref{fig:3-22} focuses on the eighteen entities introduced in this sprint, grouped into
Agile Planning and Tasks \& Kanban. The two anchor classes from earlier sprints ---
\texttt{User} (Sprint~1) and \texttt{Project} (Sprint~2) --- are shown in plain style so the
attachment points are clear: every agile artefact hangs off \texttt{Project}, and \texttt{Task}
carries optional foreign keys into \texttt{Sprint}, \texttt{Epic} and \texttt{Milestone} while
pointing back to \texttt{User} as assignee and reporter.

% Figure 3.22 — render diagrams-src/fig_3_22.puml -> figures/fig_3_22.(png|pdf)
\reportsideways{figures/fig_3_22}{Cumulative class diagram after Sprint 3.}{3-22}

\subsection{Sprint conclusion}

Sprint~3 delivered the agile planning layer and a data-driven task engine that together form the
working core of the platform: governed sprints, server-side analytics, and a Kanban board whose
columns, transitions, WIP limits, and dependency gates are enforced end-to-end. The graduated
capability model and the AI embedding outbox introduced here carry forward into every subsequent
sprint.

\fi % \ifpartialbuild (Sprints 2-3 guard opened before Sprint 2's heading above)
